\par\Needspace{7\baselineskip}\bigskip
\begin{center}
\SourceLeaf{439}
\pdfbookmark[2]{Gauss an Encke, Göttingen, 13. October 1834}{pallas-brief-18341013}
\textsc{Gauss} an \textsc{Encke}, Göttingen, 13. October 1834.
\end{center}
\noindent\hspace*{\parindent}. . . . .\quad Was endlich meine Methode der speciellen Perturbationsrechnung anlangt, so lasse ich mir gern gefallen, dass Sie solche öffentlich bekannt machen\HistoricalNoteEight{Vgl. die Aufsätze von \textsc{Encke} in den Berliner Astronomischen Jahrbüchern für 1837 und 1838, speciell die Anmerkung im erstern, S. 251.]}, da ich vorerst noch nicht weiss, wann oder ob ich selbst dazu kommen könnte und durch jenes Expedienz jedenfalls der sonst vielleicht zu besorgende indiscrete Missbrauch der Privat-Mittheilung an andere von selbst wegfällt. Meine Zahlen für die frühern Pallasoppositionen werde ich Ihnen aufsuchen, sowie die Beträge der entsprechenden Elementenstörungen: seit langer Zeit sind aber die darauf bezüglichen Papiere mir nicht durch die Hände gegangen.\quad\mbox{. . . . .}
\par
